\documentclass{beamer}
% Compilar con pdfLaTeX + BibTeX (Overleaf: opción por defecto).
\usepackage[utf8]{inputenc}
\usepackage[T1]{fontenc}
\usepackage[spanish,es-noshorthands]{babel}
\usepackage{hyperref}

% Cannot enable in Xelatex
\usepackage{pgfpages}
% \setbeameroption{hide notes} % Only slides
% \setbeameroption{show only notes} % Only notes
% \setbeameroption{show notes on second screen}

% other packages
\usepackage{latexsym,amsmath,xcolor,multicol,booktabs,calligra}
\usepackage{graphicx,listings,stackengine}
\usepackage{array}
\usepackage{natbib}

%% Enable only in Xelatex
% \usepackage{pstricks}

\author[Sierra, Carles, Miranda]{Alex Sierra \and Luis Carles \and Gonzalo Miranda Uribe}
\title{MLRanking Assignment}
\subtitle{AdaRank: training set listwise para consultas LOINC}
\institute [ETSIINF, UPM] {Escuela Técnica Superior de Ingenieros Informáticos\\Universidad Politécnica de Madrid}
\date{3 de diciembre de 2026}
\usepackage{YTU}

% defs
\def\cmd#1{\texttt{\color{red}\footnotesize $\backslash$#1}}
\def\env#1{\texttt{\color{blue}\footnotesize #1}}
\definecolor{deepblue}{rgb}{0,0,0.5}
\definecolor{deepred}{rgb}{0.6,0,0}
\definecolor{deepgreen}{rgb}{0,0.5,0}
\definecolor{halfgray}{gray}{0.55}
\newcommand{\rel}[1]{\textbf{#1}}
\newcommand{\code}[1]{\texttt{#1}}

\lstset{
    basicstyle=\ttfamily\small,
    keywordstyle=\bfseries\color{deepblue},
    emphstyle=\ttfamily\color{deepred},    % Custom highlighting style
    stringstyle=\color{deepgreen},
    numbers=left,
    numberstyle=\small\color{halfgray},
    rulesepcolor=\color{red!20!green!20!blue!20},
    frame=shadowbox,
}


\begin{document}

\begin{frame}
    \titlepage


    \begin{note}
        {Introduce your self}
    \end{note}

\end{frame}

\begin{frame}
    \tableofcontents[sectionstyle=show,subsectionstyle=show/shaded/hide,subsubsectionstyle=show/shaded/hide]
\end{frame}


% ===========================================================================
\section{Introducción}

\begin{frame}{El problema}
    \begin{itemize}
        \item Tres consultas clínicas: \emph{glucose in blood}, \emph{bilirubin in plasma}, \emph{white blood cells count}.
        \item 67 pruebas de laboratorio LOINC, \emph{las mismas} para las tres consultas: 201 pares consulta--documento.
        \item Cada documento: \code{loinc\_num}, \code{long\_common\_name}, \code{component}, \code{system}, \code{property}.
        \item El Excel no trae etiquetas: hay que \textbf{construir el training set}, ampliarlo y entrenar un ranker.
    \end{itemize}
    \vspace{6pt}
    \begin{block}{Pipeline del proyecto}
        \small Etiquetado 0--3 $\rightarrow$ ampliación (consultas y documentos) $\rightarrow$ features y LETOR
        $\rightarrow$ AdaRank + evaluación $\rightarrow$ análisis de errores
    \end{block}

    \note{Contar qué nos dieron (un Excel con 3 consultas y 67 documentos sin etiquetas) y qué había que entregar:
    un training set listwise. Presentar el pipeline: es el hilo de toda la presentación.}
\end{frame}


% ===========================================================================
\section{Estado del Arte}

% \subsection{Beamer Subsection}

\begin{frame}{Recuperación de la información (IR)}
    \begin{itemize}
        \item \textbf{Problema:} crear una función (modelo de \textit{ranking}) que pueda ordenar los documentos relevantes dada una consulta.
        \item Cuando los modelos son más sofisticados y a mayor disponibilidad de datos, el ajuste de parámetros y entrenamiento se vuelve un \textbf{desafío}.
        \item BM25, Ranking SVM, RankBoost.
    \end{itemize}

    \note {Write your notes.\\}
    \begin{note}
        {The key problem for document retrieval is ranking, specifically, how to create the ranking model (function) that can sort documents
        based on their relevance to the given query.

        It is a common practice in IR to tune the parameters of a ranking model using some labeled data and one performance measure [1]. For example, the state-of-the-art methods of BM25 [24] and LMIR (Language Models for Information Retrieval) [18, 22] all have parameters to tune. As the ranking models become more sophisticated (more features are used) and more labeled data become available, how to tune or train
        ranking models turns out to be a challenging issue.}
    \end{note}
\end{frame}

\begin{frame}{Aprendizaje automático (ML)}
    \begin{itemize}
        \item \textbf{\textit{Learning to rank}} crear una función que asigna puntuaciones a las instancias (documentos) y luego clasificarlas.
        \item \textbf{\textit{Boosting}} técnica para mejorar el rendimiento de los algoritmos de aprendizaje automático.
        \item \textbf{\textit{AdaBoost}} algoritmo que construye un modelo lineal de ‘clasificadores débiles’.
        \item \textbf{\textit{AdaRank}} método de \textit{boosting} desarrollado para ranking en IR.
    \end{itemize}

    \note {Write your notes.\\}
    \begin{note}
        {Learning to rank is to automatically create a ranking function that assigns scores to instances and then rank the instances by using the scores. Several approaches have been proposed to tackle the problem. One major approach to learning to rank is that of transforming it into binary classification on instance pairs. This ‘pair-wise’ approach fits well with information retrieval and thus is widely used in IR.

        Boosting is a general technique for improving the accuracies of machine learning algorithms. The basic idea of boosting is to repeatedly construct ‘weak learners’ by re-weighting training data and form an ensemble of weak learners such that the total performance of the ensemble is ‘boosted’.

        AdaBoost is an algorithm that ingeniously constructs a linear model by minimizing the ‘exponential loss function’ with respect to the training data [26]. Our work in this paper can be viewed as a boosting method developed for ranking, particularly for ranking in IR.}
    \end{note}
\end{frame}

\begin{frame}{Pointwise, pairwise, listwise}
    \centering\small
    \begin{tabular}{>{\bfseries}l p{3.4cm} p{3.6cm}}
        \toprule
        Enfoque & Unidad de aprendizaje & Ejemplos \\
        \midrule
        Pointwise & un documento $\rightarrow$ etiqueta & regresión, clasificación ordinal \\
        Pairwise & par $(d_i,d_j)$: ¿cuál va antes? & RankNet, Ranking SVM \\
        Listwise & la lista completa de una consulta & \alert{AdaRank}, ListNet \\
        \bottomrule
    \end{tabular}
    \vspace{8pt}
    \begin{itemize}
        \item Al usuario le importa el \emph{orden} de la lista: eso es lo que optimiza un método listwise.
        \item Por eso el training set se agrupa por consulta (\code{qid}): 3 listas de 67, no 201 filas sueltas.
    \end{itemize}

    \note{Diferencia clave: pointwise mira documentos sueltos, pairwise pares, listwise la lista entera.
    Esto justifica el formato del training set: agrupado por qid.}
\end{frame}


% ===========================================================================
\section{AdaRank}

\begin{frame}{\textit{Framework} general}
    \begin{itemize}%[<+-| alert@+>] % 当然，除了alert，手动在里面插 \pause 也行
        \item \textbf{Recuperación (prueba)} dada una consulta, el sistema retorna una lista de documentos relevantes.
        \item \textbf{Aprendizaje (entrenamiento)} construir una función que consiga los mejores resultados de clasificación al minimizar la función de pérdida.
        \item El modelo AdaRank se define como una combinación lineal de "clasificadores débiles".
        \item Utiliza medidas de desempeño de IR.
    \end{itemize}

    \note {Write your notes.\\}
    \begin{note}
        {In retrieval (testing), given a query the system returns a ranking list of documents in descending order of the relevance scores. The relevance scores are calculated with a ranking function (model). In learning (training), a number of queries and their corresponding retrieved documents are given. Furthermore, the relevance levels of the documents with respect to the queries are also provided. The relevance levels are represented as ranks (i.e., categories in a total order). The objective of learning is to construct a ranking function which achieves the best results in ranking of the training data in the sense of minimization of a loss function. Ideally the loss function is defined on the basis of the performance measure used in testing.

        The objective of learning is to create a ranking function, such that for each query the elements in its corresponding document list can be assigned relevance scores using the function and then be ranked according to the scores.

        The learning process turns out to be that of minimizing the loss function which represents the disagreement between the permutation π(qi, di, f) and the list of ranks yi, for all of the queries.}
    \end{note}
\end{frame}

\begin{frame}{Algoritmo}
    \begin{itemize}
        \item Ranker final: $f(\mathbf{x})=\sum_t \alpha_t h_t(\mathbf{x})$; cada ranker débil $h_t$ es una feature \citep{xu2007adarank}.
        \item Pesos $P_t(i)$ sobre las \textbf{consultas}, inicialmente uniformes. En cada ronda $t$:
        \begin{enumerate}
            \item elegir la feature con mejor NDCG ponderado por $P_t$;
            \item $\alpha_t=\frac12\ln\frac{\sum_i P_t(i)(1+E_i)}{\sum_i P_t(i)(1-E_i)}$;
            \item $P_{t+1}(i)\propto e^{-E(q_i,f_t)}$: las consultas mal ordenadas pesan más.
        \end{enumerate}
        \item Medida $E$: $\mathrm{NDCG@}k=\frac{1}{\mathrm{IDCG@}k}\sum_{r\le k}\frac{2^{y_r}-1}{\log_2(r+1)}$ \citep{jarvelin2002}.
    \end{itemize}

    \note {Write your notes.\\}
    \begin{note}
        {Como AdaBoost, pero las instancias son consultas. En cada ronda gana la feature que mejor ordena las consultas
        que más pesan, y las consultas que el modelo ordena mal ganan peso para la siguiente ronda.
        En la implementación corregimos un caso límite: si una feature ordena perfectamente todas las consultas, alpha es infinito.}
    \end{note}
\end{frame}

\section{Ventajas y diferencias}
\begin{frame}{Ventajas}
    \begin{itemize}
        \item Puede incorporar cualquier \textbf{medida de desempeño} de IR.
        \item El proceso de aprendizaje es más \textbf{eficiente}.
        \item Las instancias corresponden a \textbf{consultas}, no a pares de documentos.
        \item No depende del número de documentos.
    \end{itemize}

    \note {Write your notes.\\}
    \begin{note}
        {AdaRank has several other advantages when compared with the existing learning to rank methods such as Ranking SVM, RankBoost, and RankNet. First, AdaRank can incorporate any performance measure, provided that the measure is query based and in the range of [−1, +1]. Notice that the major IR measures meet this requirement. In contrast the existing methods only minimize loss functions that are loosely related to the IR measures [16].

        Second, the learning process of AdaRank is more efficient than those of the existing learning algorithms. The time complexity of
        AdaRank is of order O((k+T )·m·n log n), where k denotes the number of features, T the number of rounds, m the number of queries in training data, and n is the maximum number of documents for queries in training data.

        Third, AdaRank employs a more reasonable framework for performing the ranking task than the existing methods. Specifically in AdaRank the instances correspond to queries, while in the existing methods the instances correspond to document pairs. As a result, AdaRank does not have the following shortcomings that plague the existing methods. (a) The existing methods have to make a strong assumption that the document pairs from the same query are independently distributed. In reality, this is clearly not the case and this problem does not exist for AdaRank. (b) Ranking the most relevant documents on the tops of document lists is crucial for document retrieval. The existing methods cannot focus on the training on the tops, as indicated in [4]. Several methods for rectifying the problem have been proposed (e.g., [4]), however, they do not seem to fundamentally solve the problem. In contrast, AdaRank can naturally focus on training on the tops of document lists, because the performance measures used favor rankings for which relevant documents are on the tops. (c) In the existing methods, the numbers of document pairs vary from query to query, resulting in creating models biased toward queries with more document pairs, as pointed out in [4]. AdaRank does not have this drawback, because it treats queries rather than document pairs as basic units in learning}
    \end{note}
\end{frame}

\begin{frame}{Diferencias con \textit{AdaBoost}}
    \begin{itemize}
        \item AdaRank se desarrolló en específico para ranking en IR. AdaBoost para problemas de clasificación.
        \item AdaRank emplea medidas de desempeño de IR. AdaBoost una medida específica de clasificación.
        \item AdaRank utiliza las consultas, sus listas de documentos y niveles de relevancia. AdaBoost vectores de características y etiquetas de clase.
        \item AdaRank calcula la actualización de pesos de acuerdo al rendimiento del modelo. AdaBoost de acuerdo al rendimiento de sus "clasificadores débiles"
    \end{itemize}

    \note {Write your notes.\\}
    \begin{note}
        {First, the types of instances are different. AdaRank makes use of queries and their corresponding document lists as instances. The labels in training data are lists of ranks (relevance levels). AdaBoost makes use of feature vectors as instances. The labels in training data are simply +1 and −1. Second, the performance measures are different. In AdaRank, the performance measure is a generic measure, defined on the document list and the rank list of a query. In AdaBoost the corresponding performance measure is a specific measure for binary classification, also referred to as ‘margin’ [25]. Third, the ways of updating weights are also different. In AdaBoost, the distribution of weights on training instances is calculated according to the current distribution and the performance of the current weak learner. In AdaRank, in contrast, it is calculated according to the performance of the ranking model created so far}
    \end{note}
\end{frame}


% ===========================================================================
\section{Construcción del training set}

\begin{frame}{Escala de relevancia 0--3}
    Cada par consulta--documento se juzga en \textbf{tres dimensiones}:
    \begin{itemize}
        \item \textbf{Analito} (\code{component}): exacto / relacionado / ninguno.
        \item \textbf{Muestra} (\code{system}): exacta / compatible / incompatible.
        \item \textbf{Medida} (\code{property}, condiciones): adecuada / desviada.
    \end{itemize}
    \begin{block}{Regla (determinista)}
        \small Analito distinto $\rightarrow$ \rel{0}. Si no: base \rel{3} (exacto) o \rel{2} (relacionado);
        $-1$ muestra compatible, $-2$ incompatible; $-1$ medida desviada; mínimo \rel{1}.
    \end{block}
    \small Cada par guarda su justificación en \code{data/relevance\_labels.csv}.

    \note{La idea: no etiquetar a ojo sino con una regla que cualquiera puede repetir. Ejemplo:
    glucosa en suero/plasma para "glucose in blood": analito exacto (3), muestra compatible (-1) = 2.}
\end{frame}

\begin{frame}{Decisiones de etiquetado}
    \begin{itemize}
        \item \code{Ser/Plas} frente a \code{Bld}: compatible pero distinta $\rightarrow$ glucosa en suero/plasma = \rel{2}.
        \item \textbf{WBC = leucocitos}: linfocitos y monocitos (\#/vol) = \rel{2}; \emph{porcentaje} de linfocitos = \rel{1}.
        \item Glucosa en ayunas o post-carga: prueba funcional, medida desviada.
        \item ``Mismo término'' no es ``misma prueba'':
        \begin{itemize}
            \item glucosa en orina o en líquido pleural $\rightarrow$ \rel{1};
            \item \emph{Blood product type} para \emph{glucose in blood} $\rightarrow$ \rel{0}.
        \end{itemize}
    \end{itemize}

    \note{Estas son las decisiones que más se discutieron. El mensaje: la coincidencia de palabras
    no implica relevancia, y eso será justo lo que les cuesta a las features léxicas.}
\end{frame}

\begin{frame}{Ampliación: ¿por qué un procedimiento sistemático?}
    \begin{itemize}
        \item Primer intento: añadir términos y consultas a mano (\emph{platelets}, \emph{creatinine}).
        \item Problema: \textbf{sesgo de selección}. Casi todos los candidatos contenían la palabra de la consulta.
        \item ``Equilibrada'' \emph{no} es repartir las notas por igual: eso distorsiona el NDCG.
        \item Solución: \textbf{el mismo procedimiento para todas las consultas}, sobre las 62\,148 pruebas de LOINC.
    \end{itemize}

    \note{El piloto manual fue útil precisamente porque nos hizo ver el sesgo.}
\end{frame}

\begin{frame}{12 consultas nuevas por selección estratificada}
    \small Para cada estrato \textbf{dominio $\times$ muestra}, la prueba más frecuente (\code{COMMON\_TEST\_RANK});
    la consulta se redacta a mano.
    \vspace{4pt}

    \centering\scriptsize
    \begin{tabular}{lll}
        \toprule
        Dominio & Muestra & Consulta \\
        \midrule
        Química & suero/sangre/orina & creatinine; oxygen partial pressure; ketones \\
        Hematología & sangre & red blood cells count \\
        Coagulación & plasma & prothrombin time in plasma \\
        Microbiología & resp./orina/suero & SARS-CoV-2 RNA; urine culture; HIV Ab+Ag \\
        Uroanálisis & orina & leukocyte esterase in urine \\
        Serología & suero & perinuclear ANCA in serum \\
        Fármacos & sangre & tacrolimus in blood \\
        Patología & tejido & EGFR antigen in tissue \\
        \bottomrule
    \end{tabular}

    \note{Variedad de dominio, de muestra y de tipo de medida: concentración, recuento, tiempo, presencia...}
\end{frame}

\begin{frame}{Más documentos por \emph{pooling}}
    \begin{itemize}\small
        \item Como en TREC/LETOR, por consulta: 40 por BM25 + 20 variantes del analito + 300 aleatorios,
              con límite de casi-duplicados $\rightarrow$ \textbf{1222 documentos}.
        \item Misma regla 0--3; las 201 etiquetas oficiales no cambian.
        \item Aparecen \textbf{distractores difíciles} (relevancia 0):
    \end{itemize}
    \centering
    \resizebox{!}{0.58\textheight}{% Generado por src/export_tables.py — no editar a mano
\begin{tabular}{lp{6.6cm}lllc}
\toprule
Consulta & Documento de la colección & Analito & Muestra & Medida & Rel. \\
\midrule
oxygen partial pressure in blood & Oxygen [Partial pressure] in Blood & exacto & exacta & adecuada & 3 \\
 & Oxygen saturation in Arterial blood & relacionado & exacta & desviada & 1 \\
\midrule
prothrombin time in plasma & Prothrombin time (PT) & exacto & exacta & adecuada & 3 \\
 & INR in Platelet poor plasma by Coagulation assay & relacionado & exacta & adecuada & 2 \\
 & aPTT in Blood by Coagulation assay & ninguno & compatible & adecuada & 0 \\
\midrule
ketones in urine & Ketones [Presence] in Urine & exacto & exacta & adecuada & 3 \\
 & Methyl ethyl ketone [Mass/volume] in Urine & ninguno & exacta & adecuada & 0 \\
\midrule
SARS-CoV-2 RNA in respiratory specimen & SARS-CoV-2 (COVID-19) RNA [Presence] in Respiratory system specimen & exacto & exacta & adecuada & 3 \\
 & SARS-CoV+SARS-CoV-2 (COVID-19) Ag [Presence] in Respiratory system specimen by Rapid immunoassay & relacionado & exacta & adecuada & 2 \\
 & SARS coronavirus RNA [Presence] in Specimen by NAA with probe detection & ninguno & incompatible & adecuada & 0 \\
\midrule
bacteria in urine culture & Bacteria identified in Urine by Culture & exacto & exacta & adecuada & 3 \\
 & Bacteria identified in Bronchoalveolar lavage by Culture & exacto & incompatible & adecuada & 1 \\
\midrule
HIV antibody and antigen in serum & HIV 1+2 Ab+HIV1 p24 Ag [Presence] in Serum or Plasma by Immunoassay & exacto & exacta & adecuada & 3 \\
 & HIV 1+2 Ab [Presence] in Serum or Plasma by Immunoassay & relacionado & exacta & adecuada & 2 \\
\midrule
tacrolimus in blood & Tacrolimus [Mass/volume] in Blood & exacto & exacta & adecuada & 3 \\
 & Tacrolimus [Mass/volume] in Blood --trough & exacto & exacta & desviada & 2 \\
\midrule
EGFR antigen in tissue & Epidermal growth factor receptor Ag [Presence] in Tissue & exacto & exacta & adecuada & 3 \\
 & Epidermal growth factor [Mass/volume] in Serum or Plasma & ninguno & incompatible & desviada & 0 \\
\bottomrule
\end{tabular}
}

    \note{Ejemplos para comentar: metil etil cetona para "ketones" (es un disolvente), aPTT para "prothrombin time",
    EGF (el ligando) para EGFR (el receptor).}
\end{frame}


% ===========================================================================
\section{Features}

\begin{frame}{Seis features (sin mirar la etiqueta)}
    \centering\small
    \begin{tabular}{>{\bfseries}l p{8.6cm}}
        \toprule
        $f_1$ & coseno TF-IDF consulta--\code{long\_common\_name} \\
        $f_2$ & BM25 normalizado por consulta \citep{robertson2009} \\
        $f_3$ & solapamiento con sinónimos (\emph{white cells}$\to$\emph{leukocytes}, \emph{blood}$\to$\emph{bld}) \\
        $f_4$ & proporción del analito presente en \code{component} \\
        $f_5$ & término de muestra en \code{system} (0/1) \\
        $f_6$ & propiedad aceptada para la consulta (0/1) \\
        \bottomrule
    \end{tabular}
    \vspace{6pt}
    \begin{itemize}\small
        \item Todas dependen de la consulta. Se quitaron los \code{LabelEncoder}: son iguales para cualquier consulta.
        \item Revisión: $f_1$ y $f_2$ ordenan casi igual (Spearman $>0.99$).
    \end{itemize}

    \note{Una feature solo sirve como ranker débil si cambia según la consulta. f2 resultó redundante con f1.}
\end{frame}

\begin{frame}{Ajuste por popularidad del código LOINC}
    \begin{itemize}
        \item Idea: a igualdad de texto, la prueba de uso común es la que busca el clínico.
        \item Popularidad a partir de \code{COMMON\_TEST\_RANK} ($r_d$, menor = más usada):
        \[ P(d)=1-\frac{\log(1+r_d)}{\log(1+20\,000)} \in [0,1] \]
        \item Ajuste suave de $f_1$ y $f_3$: \quad $f_i' = f_i\,[1+0.1\,(2P(d)-1)]$, \quad multiplicador entre 0.9 y 1.1.
        \item Como feature independiente \emph{no} mejoraba; como ajuste, sí.
    \end{itemize}

    \note{No cambia el orden entre documentos con texto muy distinto; sobre todo desempata entre documentos parecidos.}
\end{frame}

\begin{frame}[fragile]{Formato LETOR}
    \small
    \begin{itemize}
        \item \code{train.txt} (3$\times$67, obligatorio), \code{train\_terms.txt} (3$\times$1222), \code{train\_extended.txt} (15$\times$1222).
        \item Una línea por par: etiqueta, \code{qid}, features \code{índice:valor} y el código LOINC en comentario.
    \end{itemize}
    % Generado por src/export_tables.py — no editar a mano
\begin{lstlisting}
3 qid:1 1:0.750959 2:1 3:1.01765 4:1 5:1 6:1 # docid = 74774-1
2 qid:1 1:0.499538 2:0.694848 3:0.465047 4:1 5:0 6:1 # docid = 14749-6
1 qid:1 1:0.428328 2:0.75537 3:0.5 4:1 5:0 6:1 # docid = 15076-3
3 qid:2 1:0.653337 2:1 3:1.03668 4:1 5:1 6:1 # docid = 1975-2
3 qid:3 1:0.467694 2:0.965348 3:1.01449 4:0.333333 5:1 6:1 # docid = 26464-8
1 qid:3 1:0.308181 2:0.881265 3:0.761091 4:0.333333 5:1 6:0 # docid = 26478-8
0 qid:1 1:0.237446 2:0.458793 3:0.491787 4:0 5:0 6:0 # docid = 933-2
\end{lstlisting}


    \note{Es el formato estándar de LETOR/RankLib; nuestro AdaRank lo lee directamente.}
\end{frame}


% ===========================================================================
\section{Resultados y discusiones}

\begin{frame}{Evaluación leave-one-query-out}
    \small Cada fold entrena con todas las consultas menos una y evalúa en la restante (NDCG medio).
    \vspace{4pt}

    \centering
    \resizebox{\linewidth}{!}{% Generado por src/export_tables.py — no editar a mano
\begin{tabular}{lrrrr}
\toprule
Sistema & NDCG@1 & NDCG@3 & NDCG@5 & NDCG@10 \\
\midrule
\multicolumn{5}{l}{\emph{3 consultas oficiales (3 folds)}} \\
\textbf{AdaRank} & 0.667 & 0.772 & 0.796 & 0.797 \\
solo $f_1$ (coseno TF-IDF) & 0.667 & 0.856 & 0.889 & 0.889 \\
solo $f_2$ (BM25) & 0.667 & 0.764 & 0.829 & 0.863 \\
solo $f_3$ (solap. sinónimos) & 1.000 & 0.875 & 0.885 & 0.877 \\
solo $f_4$ (analito en component) & 0.714 & 0.824 & 0.821 & 0.824 \\
solo $f_5$ (muestra en system) & 0.333 & 0.366 & 0.475 & 0.467 \\
\midrule
\multicolumn{5}{l}{\emph{6 consultas, ampliado$^\dagger$ (6 folds)}} \\
\textbf{AdaRank} & 0.833 & 0.928 & 0.944 & 0.944 \\
solo $f_1$ (coseno TF-IDF) & 0.833 & 0.928 & 0.944 & 0.944 \\
solo $f_2$ (BM25) & 0.833 & 0.882 & 0.914 & 0.931 \\
solo $f_3$ (solap. sinónimos) & 0.905 & 0.878 & 0.878 & 0.893 \\
solo $f_4$ (analito en component) & 0.619 & 0.820 & 0.832 & 0.834 \\
solo $f_5$ (muestra en system) & 0.167 & 0.183 & 0.237 & 0.234 \\
\bottomrule
\end{tabular}
}

    \note{Con 3 o 15 consultas no hay partición train/test fiable; por eso leave-one-query-out.
    Cada feature sola es una línea base.}
\end{frame}

\begin{frame}{Efecto del ajuste por popularidad}
    \centering\small
    \begin{tabular}{lcc}
        \toprule
        NDCG@10 de AdaRank & Antes & Con popularidad \\
        \midrule
        Oficial (3$\times$67) & 0.798 & \textbf{0.980} \\
        + términos (3$\times$1222) & 0.433 & \textbf{0.482} \\
        + términos y consultas (15$\times$1222) & 0.627 & \textbf{0.695} \\
        \bottomrule
    \end{tabular}
    \vspace{8pt}
    \begin{itemize}\small
        \item \emph{white blood cells count} en la colección completa: el recuento total de leucocitos pasa del puesto 21 al \textbf{1}.
        \item AdaRank acaba eligiendo \textbf{una sola feature} ($f_1$ en el oficial, $f_3$ en la colección completa).
    \end{itemize}

    \note{"Antes" es la versión anterior de las features (sin popularidad). La mejora es clara, pero ojo:
    las consultas nuevas se eligieron como la prueba más frecuente de su estrato, así que parte de la mejora puede venir de ahí.}
\end{frame}

\begin{frame}{¿Dónde falla? Análisis de errores}
    \begin{itemize}\small
        \item Va bien si el nombre LOINC contiene la consulta: \emph{leukocyte esterase} 0.96, \emph{urine culture} 0.90, \emph{prothrombin time} 0.88.
        \item Falla si no: \emph{red blood cells count} 0.49, \emph{SARS-CoV-2 RNA} 0.51.
        \item Falsos positivos (poco relevantes en el top-10): fracciones y cocientes, p.\,ej. \emph{Bilirubin.direct/Bilirubin.total}.
        \item Falsos negativos (relevantes por debajo del 10): muestras especiales, p.\,ej. \emph{Glucose in Capillary blood}.
    \end{itemize}
    \begin{block}{Conclusión}
        \small Las features léxicas no ven los matices que separan un \rel{3} de un \rel{1} (muestra, condición, cociente).
    \end{block}

    \note{El análisis de errores cruza el ranking con las etiquetas y nos dice qué tipo de documento se cuela o se pierde.
    Es lo que orienta qué features añadir.}
\end{frame}


% ===========================================================================
\section{Conclusiones}

\begin{frame}{Conclusiones y limitaciones}
    \begin{itemize}
        \item Training set listwise \textbf{reproducible}: regla 0--3 documentada, ampliación sistemática, 6 features, LETOR y comprobaciones.
        \item La popularidad es la mejora más clara; la variedad de consultas revela dónde fallan las features.
        \item Limitaciones:
        \begin{itemize}
            \item 15 consultas siguen siendo pocas: resultados demostrativos;
            \item posible sesgo de la popularidad por cómo se eligieron las consultas;
            \item $f_4$--$f_6$ usan las mismas columnas que la regla de etiquetado.
        \end{itemize}
        \item Siguiente paso: features semánticas o estructurales del término LOINC.
    \end{itemize}
\end{frame}

\begin{frame}{Reparto del trabajo}
    \centering\small
    \begin{tabular}{>{\bfseries}p{2.6cm} p{7.4cm}}
        \toprule
        Gonzalo & Implementación base de AdaRank y métricas, primer dataset (0--2) y notebook;
                  ampliación piloto; estudio de AdaRank y presentación. \\
        \addlinespace
        Alex & Escala 0--3 y guía de etiquetado; selección estratificada y pooling;
               pipeline reproducible, LETOR y documentación. \\
        \addlinespace
        Luis & Revisión de features (Spearman, $f_4$, $f_6$); ajuste por popularidad;
               análisis de errores. \\
        \bottomrule
    \end{tabular}
\end{frame}


% ===========================================================================
\section{Referencias}
\begin{frame}[allowframebreaks]
    \nocite{chen-adarank}
    \bibliography{ref}
    \bibliographystyle{alpha}
    % If too many references, use this command to resize：
    %\tiny\bibliographystyle{alpha}
\end{frame}

\begin{frame}
    \begin{center}
        {\Huge ¡Gracias!}

        \vspace{12pt}
        {\large ¿Preguntas?}
    \end{center}
\end{frame}

\end{document}
